% =============================================================================
% Proposta de Trabalho - Disciplina INF0330 (UFG)
% Sistema de Gestão de Biblioteca com Recomendações
%
% Este arquivo segue configuração ABNT/UFG básica (Times 12pt, margens 3/2 cm,
% espaçamento 1,5). Compila standalone com pdflatex no Overleaf.
%
% Se quiser usar o template oficial Tese UFG do Overleaf, copie o conteúdo
% entre \begin{document} e \end{document} para lá, ajustando as seções
% conforme a estrutura do template (frontmatter, mainmatter, etc.).
% =============================================================================

\documentclass[12pt,a4paper]{article}

% Codificação e idioma
\usepackage[utf8]{inputenc}
\usepackage[T1]{fontenc}
\usepackage[brazilian]{babel}

% Fonte Times Roman (padrão ABNT)
\usepackage{mathptmx}

% Margens ABNT (3 cm superior/esquerda, 2 cm inferior/direita)
\usepackage[top=3cm, left=3cm, right=2cm, bottom=2cm]{geometry}

% Espaçamento 1,5 entrelinhas
\usepackage{setspace}
\onehalfspacing

% Indentação do primeiro parágrafo
\usepackage{indentfirst}
\setlength{\parindent}{1.25cm}

% Tabelas e listas
\usepackage{booktabs}
\usepackage{array}
\usepackage{tabularx}
\usepackage{enumitem}

% Links
\usepackage[hidelinks]{hyperref}

% Formatação de seções (caixa-alta no nível \section)
\usepackage{titlesec}
\titleformat{\section}{\normalfont\bfseries\large\MakeUppercase}{\thesection}{1em}{}
\titleformat{\subsection}{\normalfont\bfseries\normalsize}{\thesubsection}{1em}{}

% Sem indentação em itens de lista
\setlist[itemize]{leftmargin=*}

% =============================================================================
\begin{document}

% -------------------- Folha de rosto --------------------
\begin{titlepage}
\centering
\vspace*{2cm}

{\large UNIVERSIDADE FEDERAL DE GOIÁS}\\[0.3cm]
{\large PÓS-GRADUAÇÃO}\\[0.3cm]
{\large DISCIPLINA INF0330 --- FRAMEWORK DE DESENVOLVIMENTO WEB PARA CONSUMO DE MODELOS TREINADOS DE INTELIGÊNCIA ARTIFICIAL}\\[3cm]

{\Large \textbf{SISTEMA DE GESTÃO DE BIBLIOTECA COM RECOMENDAÇÕES}}\\[0.5cm]
{\large Proposta de Trabalho da Disciplina}\\[3cm]

\begin{flushright}
\begin{minipage}{0.55\textwidth}
\raggedright
\textbf{Autores:}\\
Antonio Jansen\\
Givanildo Gramacho\\
Jucelino Santos\\
Ronny Marcelo\\
Vanderson Darriba\\[0.8cm]

\textbf{Professor:} Ronaldo M. da Costa
\end{minipage}
\end{flushright}

\vfill
{\large Goiânia, abril de 2026}

\end{titlepage}

% -------------------- Seção 1 --------------------
\section{Componentes do Grupo}

\begin{itemize}
    \item Antonio Jansen
    \item Givanildo Gramacho
    \item Jucelino Santos
    \item Ronny Marcelo
    \item Vanderson Darriba
\end{itemize}

% -------------------- Seção 2 --------------------
\section{Contextualização}

A disciplina INF0330 tem como foco o desenvolvimento de aplicações web utilizando frameworks modernos, com ênfase na integração e consumo de modelos treinados de Inteligência Artificial em cenários reais de uso. O trabalho proposto neste documento busca atender diretamente essa ementa, combinando um sistema web tradicional (CRUD com autenticação, gestão de entidades e regras de negócio) com uma camada inteligente de recomendação baseada em modelo de linguagem treinado.

% -------------------- Seção 3 --------------------
\section{Definição do Trabalho}

\subsection{Problema}

Bibliotecas acadêmicas lidam diariamente com grandes acervos de livros, teses, dissertações e monografias. Os leitores (alunos, pesquisadores e docentes) enfrentam duas dificuldades recorrentes:

\begin{enumerate}
    \item A gestão operacional do acervo é trabalhosa sem um sistema digital adequado (cadastros, empréstimos, devoluções, controle de atrasos).
    \item A descoberta de novas obras relevantes dentro do acervo é limitada, uma vez que os mecanismos tradicionais de busca dependem de correspondência exata por título ou autor, deixando de explorar relações semânticas entre os conteúdos.
\end{enumerate}

\subsection{Solução Proposta}

Desenvolver um sistema web de gestão de biblioteca que cubra o ciclo completo de operação do acervo e que, adicionalmente, ofereça recomendações de leituras baseadas em similaridade semântica. O sistema utilizará um modelo treinado de representação textual para gerar sugestões personalizadas ao leitor, a partir do livro consultado e do seu histórico de empréstimos.

\subsection{Objetivos}

\noindent\textbf{Objetivo geral.} Construir uma aplicação web funcional e segura, desenvolvida com o framework Django, que integre um módulo de recomendação baseado em modelo treinado de Inteligência Artificial.

\noindent\textbf{Objetivos específicos.}

\begin{itemize}
    \item Modelar e implementar o CRUD das entidades da biblioteca (pessoas, obras, empréstimos).
    \item Implementar regras de negócio de empréstimo, devolução, controle de disponibilidade e atrasos.
    \item Prover autenticação e autorização por perfis de usuário (bibliotecário, leitor e administrador).
    \item Consumir um modelo treinado para gerar recomendações de obras por similaridade semântica.
    \item Disponibilizar um painel de indicadores (\textit{dashboard}) com métricas operacionais do acervo.
\end{itemize}

% -------------------- Seção 4 --------------------
\section{Escopo do Sistema}

\subsection{Módulo CRUD (base operacional)}

\noindent\textbf{Entidades.} Pessoa (bibliotecário e leitor), Livro (bibliografia, tese/dissertação, monografia) e Empréstimo.

\noindent\textbf{Funcionalidades.}

\begin{itemize}
    \item Cadastro, edição, consulta e remoção das três entidades.
    \item Autenticação com \textit{login} e grupos de permissão (Editor, Visualizador, Superusuário).
    \item Regras de negócio no empréstimo: validação de disponibilidade do exemplar, validação de leitor ativo, cálculo automático da data de devolução prevista (14 dias) e atualização do acervo ao registrar devolução.
    \item Cálculo de \textit{status} dinâmico do empréstimo (Emprestado, Devolvido, Atrasado).
    \item Relatório de empréstimos atrasados.
    \item \textit{Dashboard} com contadores e indicadores básicos.
\end{itemize}

\subsection{Módulo de Inteligência Artificial}

\noindent\textbf{Funcionalidade principal: Recomendação por Similaridade Semântica.}

\begin{itemize}
    \item Gerar representações vetoriais (\textit{embeddings}) dos títulos e sinopses das obras do acervo.
    \item Ao visualizar uma obra, apresentar ao leitor uma lista de obras similares do acervo com base na distância vetorial (similaridade cosseno).
    \item Considerar o histórico de empréstimos do leitor para personalizar as sugestões.
\end{itemize}

\noindent\textbf{Modelos candidatos em avaliação.} \textit{Sentence-Transformers} (HuggingFace) para geração de \textit{embeddings} multilíngues. Modelos de linguagem abertos como Llama ou DeepSeek poderão ser avaliados para uma extensão conversacional. A escolha final será registrada após uma prova de conceito durante a primeira \textit{sprint} da fase de IA.

\noindent\textbf{Extensão opcional.} Caso o cronograma permita, será implementado um assistente conversacional sobre o acervo, capaz de responder perguntas em linguagem natural a respeito da disponibilidade e do conteúdo das obras cadastradas.

\subsection{Fora do Escopo}

\begin{itemize}
    \item Aplicação móvel nativa.
    \item Integração com sistemas de pagamento ou cobrança automática de multas.
    \item Digitalização do conteúdo integral das obras (somente metadados e sinopses).
    \item Publicação em ambiente de produção com infraestrutura de nuvem gerenciada.
\end{itemize}

% -------------------- Seção 5 --------------------
\section{Stack Tecnológica}

\begin{itemize}
    \item \textbf{Linguagem:} Python 3.x
    \item \textbf{Framework web:} Django 4.2
    \item \textbf{Banco de dados:} SQLite (desenvolvimento)
    \item \textbf{Interface:} Bootstrap 5 via \textit{django-bootstrap-v5} e \textit{django-tables2}
    \item \textbf{Autenticação:} sistema nativo do Django com grupos e permissões
    \item \textbf{Camada de IA:} biblioteca \textit{Sentence-Transformers} (HuggingFace); avaliação de modelos abertos (Llama, DeepSeek) para extensão conversacional
    \item \textbf{Controle de versão:} Git e GitHub
    \item \textbf{Ambiente de apresentação:} servidor de desenvolvimento local
\end{itemize}

% -------------------- Seção 6 --------------------
\section{Cronograma}

O trabalho está organizado em quatro \textit{sprints}, totalizando aproximadamente sete semanas de execução, com entrega final prevista para 06 de junho de 2026.

\vspace{0.3cm}

\begin{center}
\begin{tabularx}{\textwidth}{@{} l l X @{}}
\toprule
\textbf{Sprint} & \textbf{Período} & \textbf{Entregas} \\
\midrule
Sprint 0 & 21/04 a 26/04 & Documento de proposta submetido \\
Sprint 1 & 27/04 a 10/05 & Prova de conceito do módulo de recomendação; definição final do modelo \\
Sprint 2 & 11/05 a 24/05 & Integração do módulo de recomendação ao CRUD; interface de sugestões \\
Sprint 3 & 25/05 a 06/06 & Testes, \textit{dashboard} de métricas, revisão de segurança, documentação final e apresentação \\
\bottomrule
\end{tabularx}
\end{center}

\vspace{0.3cm}

Revisões semanais em grupo serão realizadas para acompanhamento de \textit{backlog}, alinhamento técnico e ajustes de escopo, se necessário.

% -------------------- Seção 7 --------------------
\section{Divisão de Responsabilidades}

Todos os integrantes participam das decisões arquiteturais, das revisões de código e das retrospectivas de \textit{sprint}. Os papéis abaixo indicam a liderança de cada frente, e não exclusividade.

\vspace{0.3cm}

\begin{center}
\begin{tabularx}{\textwidth}{@{} l X @{}}
\toprule
\textbf{Integrante} & \textbf{Frente Principal} \\
\midrule
Antonio Jansen & Arquitetura do sistema, coordenação, redação e integração da camada de IA \\
Jucelino Santos & \textit{Backend}, modelagem de dados, integração do motor de recomendação \\
Givanildo Gramacho & Pesquisa, prova de conceito e avaliação dos modelos de IA \\
Vanderson Darriba & Segurança, autenticação, conformidade e boas práticas \\
Ronny Marcelo & Métricas, análise quantitativa e painel de indicadores \\
\bottomrule
\end{tabularx}
\end{center}

% -------------------- Seção 8 --------------------
\section{Resultados Esperados}

Ao final da disciplina, o grupo entregará:

\begin{itemize}
    \item Uma aplicação web funcional para gestão de biblioteca, com autenticação, CRUD completo e regras de negócio implementadas.
    \item Um módulo de recomendação por similaridade semântica integrado ao sistema, consumindo um modelo treinado de linguagem.
    \item Documentação técnica do sistema, instruções de execução e relatório de avaliação da solução.
    \item Apresentação demonstrando o funcionamento \textit{end-to-end} da aplicação.
\end{itemize}

\end{document}
